\FC{hardware-club-001.mp3}{\HCtitle{\#1}{Tak nám ten rok pěkně začal}}{2018-02-01}{\Sleepless, \Janchor, \JSvak, \JZubec}{}{(02:58) Spectre/Meltdown | (21:16) Nedostatek grafických karet na trhu | (35:43) Notebook s GTX 1080 Max-Q (Asus Zephyrus) | (54:42) CES 2018 a GeForce NOW}{1:03:52}{41250060}{}
\FC{hardware-club-002.mp3}{\HCtitle{\#2}{Procesory}}{2018-02-12}{\Sleepless, \Janchor, \JSvak, \JZubec}{}{(02:04) Jak jsou na tom dnes procesory? (Ryzen/Coffee Lake) | (16:37) Přetaktování CPU | (20:00) Upgrade, nebo počkat? | (30:24) Operační paměti}{0:55:03}{36100368}{}
\FC{hardware-club-003.mp3}{\HCtitle{\#3}{AMD Raven Ridge}}{2018-02-22}{\Sleepless, \Janchor, \JSvak, \JZubec}{}{(03:02) Přestupy Intel/AMD | (10:58) Trh GPU a AMD | (20:11( APU AMD (Raven Ridge) | (45:09) Cenové srovnání: APU vs. CPU+GPU}{1:00:16}{40688460}{}
\FC{hardware-club-004.mp3}{\HCtitle{\#4}{Přetaktování}}{2018-03-09}{\Sleepless, \Janchor, \JSvak, \JZubec}{}{(03:25) Přetaktování | (09:56) Odbočka – frekvence vs. jádra | (15:22) Na čem jde přetaktovat a jak to funguje?  | (24:44) Náklady na OC I. | (32:50) Odbočka: 4K + VR | (46:43) Náklady na OC II. + podvoltování | (56:48) Dotazy}{1:11:56}{29674152}{}
% \FC{hardware-club-005.mp3}{\HCtitle{\#5}{Výrobní procesy, taktování AMD, delid}}{2018-03-22}{\Sleepless, \Janchor, \JSvak, \JZubec}{}{(01:17) Konec křemíku: 7nm výrobní proces a dále | (14:08) VSUVKA: Kvantové počítače | (20:37) Výrobní procesy – pokračování, pozice Intelu | (33:46) Taktování AMD, výhled s novou generací | (46:15) Delid procesoru | (56:52) Dotazy}{}{}{}
\FC{hardware-club-006.mp3}{\HCtitle{\#6}{Ray-tracing, notebook HP Omen X}}{2018-04-05}{\Sleepless, \Janchor, \JSvak, \JZubec}{}{Ray-tracing – v herní branži periodicky se objevující termín dostal novou krev na letošním GDC. Nvidia a další technologičtí giganti oznámili, že se na tento „svatý grál“ vykreslování hodlají v blízké budoucnosti zaměřit. Společně si toto téma probereme. V druhé půli se koukneme na momentálně recenzovaný notebook HP Omen X s GTX 1080. Jak vypadá v porovnání s dříve testovaným Asus ROG Zephyrus?}{1:11:03}{46254580}{}
\FC{hardware-club-007.mp3}{\HCtitle{\#7}{Ryzen 2700X \& 2600X, osmijádro od Intelu a pokažené Atomy}}{2018-04-19}{\Janchor, \JSvak, \JZubec}{}{(01:11) Druhá generace procesorů Ryzen | (14:31) Nově objevené chyby platformy Intel | (27:16) Klepy o osmijádru Intel | (31:25) Alternativa z diskuze: starší serverový Xeon? | (38:19) Ryzen 7 2700X | (48:13) Dotazy}{1:10:40}{47498688}{}
\FC{hardware-club-008.mp3}{\HCtitle{\#8}{Grafické karty}}{2018-05-03}{\Sleepless, \Janchor, \JSvak, \JZubec}{}{Trh s GPU regeneruje - Kdy jít do koupě? - Kvantum dotazů}{1:13:23}{48592152}{}
\FC{hardware-club-009.mp3}{\HCtitle{\#9}{Kdy vyjdou nové herní konzole a další novinky}}{2018-05-18}{\Sleepless, \Janchor, \JSvak, \JZubec}{}{(01:40) Novinky | (03:14) Nvidia Turing již v červenci? | (14:50) AMD Vega 20 a 7nm proces | (20:05) Co by měl vlastně přinést 7nm proces? | (26:26) Herní konzole – kdy by mohly reálně vyjít? | (50:18) Dotazy}{1:05:34}{28564504}{}
\FC{hardware-club-010.mp3}{\HCtitle{\#10}{Monitory}}{2018-05-31}{\Sleepless, \Janchor, \JSvak, \JZubec}{}{(00:57) Téma: monitory | (02:58) CRT? | (03:49) Hlas mimo kameru | (04:33) LCD TV vs. LED TV | (05:49) Typy LCD panelů | (06:43) Odezva | (12:28) Jindrův favorit mezi monitory | (19:08) IPS vs. TN | (22:54) VA panel, pozorovací úhly | (25:37) Srovnání panelů - pokračování | (34:09) Obnovovací frekvence | (43:29) GSync/FreeSync | (50:10) Dotaz: levné monitory | (54:52) Dotaz: OLED/MicroLED}{1:04:22}{24719436}{}
\FC{hardware-club-011.mp3}{\HCtitle{\#11}{Monitory - HDR, ultrawide a naše prognózy}}{2018-06-21}{\Janchor, \JSvak, \JZubec}{}{(01:33) Monitory podruhé | (03:25) Mini LED | (10:41) Co je to HDR? | (11:30) HDR: foto vs. video | (14:03) Jak HDR funguje? | (17:17) Monitor 4:3 HDR? | (21:04) Dotaz: HDR settings ve hrách | (22:46) Podmínky HDR – jas, gamut, barevná hloubka, HDR formáty | (26:28) Režie HDR obrazu | (29:35) Podpora HDR | (38:50) Ultrawide | (47:51) 4K HDR 144 Hz G-Sync displeje a jejich problémy | (56:25) Dotazy | (1:09:52) Jak budou vypadat monitory za 10 let? | (1:19:08) Jindrova sestava | (1:20:40) [soupis sestavy]}{1:24:59}{32265144}{}
\FC{hardware-club-012.mp3}{\HCtitle{\#12}{Navrhujeme herní PC do 15\,000, 30\,000 až 50\,000+ Kč}}{2018-07-03}{\Sleepless, \Janchor, \JSvak, \JZubec}{}{(03:19) Herní PC za 15\,000\,Kč | (23:52) Sestava za 30\,000\,Kč | (43:52) High-end kolem 50\,000\,Kč | (47:54) Herní dělo pro přetaktování | (53:36) Workstation za 100\,000\,Kč}{1:12:47}{31475068}{}
\FC{hardware-club-013.mp3}{\HCtitle{\#13}{Virtuální realita s Markem Ďurišem}}{2018-07-19}{\Janchor, \JSvak, \MDuris}{}{(05:02) Představení VR headsetů a naše zkušenosti | (11:28) Hardwarové požadavky na VR (PC) | (17:12) Sledování pozice | (19:59) Rozdíly v headsetech? | (22:18) Prodeje VR | (24:18) Kam s kabely? | (27:57) Bezdrátové VR | (29:58) Nevýhody a kompromisy VR | (43:25) VR a obsah pro dospělé | (45:18) VR pro mobily | (47:53) Killer hra pro VR? | (51:25) Dotazy | (1:01:15) Budoucnost VR}{1:08:47}{29154306}{}
\FC{hardware-club-014.mp3}{\HCtitle{\#14}{Chlazení}}{2018-08-02}{\Sleepless, \JSvak, \JZubec}{}{(02:09) Airflow | (06:58) Konstrukce PC skříně | (11:46) Vzduchové chlazení | (18:18) Chladič CPU + jak fungují heatpipes? | (26:14) Chlazení GPU | (35:09) Vodní chlazení | (53:32) Chlazení olejem nebo pasivně? | (58:15) Tipy na chladiče | (1:20:36) Dotazy}{1:28:58}{39573786}{}
\FC{hardware-club-015.mp3}{\HCtitle{\#15}{Herní notebooky}}{2018-08-16}{\Sleepless, \Janchor, \JSvak, \JZubec}{}{(04:24) Proč si lidé kupují herní notebooky? | (10:58) Nevýhody herních notebooků | (14:25) Externí GPU? | (16:08) Co už je herní notebook? A ceny? | (25:55) Nepřehlednost nabídky | (33:57) Poměr cena/výkon desktop vs. notebook | (38:20) Undervolting | (46:12) Modulární konfigurace | (51:15) Dotazy}{1:11:21}{30678093}{}
\FC{hardware-club-015-spec.mp3}{\HCtitle{SPECIÁL}{Nvidia Geforce RTX (Gamescom 2018)}}{2018-08-23}{\Janchor, \JSvak, \JZubec}{}{(03:04) RTX 2080Ti, 2080, 2070 | (04:55) Ray tracing a Turing | (16:18) Ceny | (20:52) Úvaha nad strategií Nvidia | (26:59) Prosadí se ray tracing? | (52:05) Parametry čipů Turing}{1:06:49}{29474397}{}
\FC{hardware-club-016.mp3}{\HCtitle{\#16}{Zdražování}}{2018-09-06}{\Janchor, \JSvak, \JZubec}{}{ (02:57) Drahé novinky ze stáje Nvidia | (15:54) Intel a nebezpečí nedostatku výrobních kapacit | (24:55) i9-9900K / i7-9700K | (28:54) Global Foundries a zastavený vývoj 7nm výrobního procesu | (32:37) BFGD | (38:17) Příležitost pro AMD? | (41:24) Stav na trhu s RAM | (45:09) Nová konzole v Číně | (47:20) AMD GPU na 7nm | (48:50) Prolomení bobříka mlčení | (50:29) Dotazy}{1:03:01}{27692316}{}
\FC{hardware-club-017.mp3}{\HCtitle{\#17}{GeForce RTX – shrnutí recenzí a naše dojmy}}{2018-09-20}{\Janchor, \JSvak, \JZubec}{}{Kde je ray tracing? | Konečně grafika pro 4K / 60 FPS? | GeForce RTX 2080 od Gigabyte}{0:54:31}{23630013}{}
\FC{hardware-club-018.mp3}{\HCtitle{\#18}{SSD disky}}{2018-10-04}{\Janchor, \JSvak, \JZubec}{}{(02:50) Co je to SSD? | (09:34) Životnost | (12:49) Rozhraní | (23:25) Paměťové čipy | (46:48) Spolehlivost | (55:00) Kapacita vs. rychlost}{1:14:31}{32863401}{}
\FC{hardware-club-019.mp3}{\HCtitle{\#19}{Nvidia RTX 2070, RAID}}{2018-10-18}{\Janchor, \JSvak, \JZubec}{}{(02:32) Doháníme dotazy | (19:06) RTX 2070: výkon, cena | (34:49) RAID | (48:19) RAID jako záloha? | (50:27) RAID 5, 6, 10 | (55:02) Jiné řešení redundance}{1:06:22}{29451888}{}
\FC{hardware-club-020.mp3}{\HCtitle{\#20}{Intel Core i9-9900K a i7-9700K}}{2018-11-01}{\Sleepless, \JSvak, \JZubec}{}{(01:52) i9-9900K a i7-9700K - parametry a výkon | (19:17) Provozní vlastnosti nových osmijader | (29:33) TDP a vliv základních desek | (39:18) Vyplatí se nové procesory? | (55:50) Opravy Meltdown/Spectre | (1:03:03) Dotazy}{1:23:26}{37838925}{}
\FC{hardware-club-021.mp3}{\HCtitle{\#21}{Externí GPU}}{2018-11-15}{\Janchor, \JSvak, \JZubec}{}{(02:36) Co je to externí grafika? | (06:48) Jak vyřešit připojení externí GPU? | (14:10) Existuje prostor pro eGPU? Výhody a nevýhody | (38:46) Novinka: RTX v Battlefield V | (59:40) Nadějný výhled AMD | (1:00:40) Dotazy | (1:07:21) Upoutávka na příště}{1:08:39}{31751073}{}
\FC{hardware-club-022.mp3}{\HCtitle{\#22}{Navrhujeme herní počítače pod stromeček}}{2018-11-29}{\Janchor, \JSvak, \JZubec}{}{(04:14) Sestava \uv{z nouze ctnost} (8\,450\,Kč) | (15:38) Sestava \uv{i Skrblík by se zaradoval} (16\,600--17\,500\,Kč) | (28:10) Sestava MID aka střední třída (23\,700--30\,900\,Kč) | (42:36) Sestava HIGH aka nejvyšší třída včetně OC (56\,600--66\,600\,Kč)}{1:05:02}{30793392}{}
\FC{hardware-club-023.mp3}{\HCtitle{\#23}{AMD a Intel v roce 2019, konkurence ARMu a další novinky}}{2018-12-13}{\Janchor, \JSvak, \JZubec}{}{(16:49) Drby o procesorech AMD Ryzen 3000 | (38:49) Nové Radeony příští rok? | (47:47) Intel a nová 10nm jádra Sunny Cove | (57:17) Ohrožení Intelu procesory ARM | (1:07:06) Dotazy / očekávání roku 2019}{1:13:12}{34443324}{}
\FC{hardware-club-024.mp3}{\HCtitle{\#24}{To nejdůležitější z CES 2019}}{2019-01-14}{\Sleepless, \Janchor, \JSvak, \JZubec}{}{(02:30) Nvidia GeForce RTX 2060: výkon a cena | (10:39) FreeSync na Nvidia kartách? | (26:17) AMD na CES: prezentace Lisy Su | (31:47) AMD a nová konfigurace CPU | (34:10) Jak budou vypadat nové Ryzeny? | (39:49) Stane se z AMD druhý Intel? | (49:23) Nová grafika AMD Radeon VII: výkon, cena, dostupnost | (54:01) Dotazy I | (1:04:09) Intel na závěr | (1:09:14) Dotazy II}{1:13:11}{36866961}{}
\FC{hardware-club-025.mp3}{\HCtitle{\#25}{Mini PC}}{2019-01-24}{\Sleepless, \Janchor, \JSvak, \JZubec}{}{(05:33) Co je Mini PC | (21:48) Chlazení | (27:25) Grafika vs iGPU | (37:25) Pasivní HTPC | (47:45) Intel NUC | (50:17) AMD APU}{1:07:56}{32853753}{}
\FC{hardware-club-026.mp3}{\HCtitle{\#26}{AI neboli umělá inteligence}}{2019-02-14}{\Janchor, \JSvak, \JZubec}{}{(01:53) Průmyslová revoluce v4.0? | (04:46) Vývoj AI | (09:55) Hardware AI | (24:45) Princip AI | (1:03:30) Současné aplikace AI a potenciální problémy | (1:54:27) AI jako další vývojový stupeň člověka?}{2:08:21}{62102196}{}
\FC{hardware-club-027.mp3}{\HCtitle{\#27}{AI neboli umělá inteligence – pokračování}}{2019-02-28}{\Janchor, \JSvak, \JZubec}{}{(01:02) Budoucnost je dnes! | (05:36) Osud lidstva v éře AI? | (16:00) Superinteligence | (23:54) Co náš čeká v nejbližších letech? | (33:13) Pozitivní/negativní vize světa s AI | (49:25) AI ve sci-fi dílech}{1:28:53}{42222816}{}
\FC{hardware-club-028.mp3}{\HCtitle{\#28}{Vyplatí se dnes hotové sestavy z obchodu?}}{2019-03-14}{\Janchor, \JSvak, \JZubec}{}{(05:13) Co se změnilo na trhu s hotovými PC sestavami? | (09:09) Srovnání nejprodávanějších PC sestav v ČR | (26:26) Výhody a nevýhody hotových sestav | (41:47) Kvalita sestavení | (54:31) Konfigurovatelná PC | (57:56) Rady na závěr}{1:08:14}{32157936}{}
\FC{hardware-club-029.mp3}{\HCtitle{\#29}{Ray tracing pro všechny, výkon GTX 1660 / 1660 Ti a další novinky}}{2019-03-28}{\Janchor, \JSvak, \JZubec}{}{(01:23) Ray tracing bez karet Nvidia RTX | (13:51) GTX 1660 / 1660 Ti: výkon ve Full HD | (18:45) Poměr cena/výkon | (25:06) Výkon ve QHD | (32:04) Spotřeba | (33:57) Ryzen 3000: 16/32 jader!}{0:52:01}{24322909}{}
\FC{hardware-club-030.mp3}{\HCtitle{\#30}{Kolik paměti je potřeba v herním PC?}}{2019-04-15}{\Janchor, \JSvak, \JZubec}{}{(02:20) RAM vs VRAM | (13:14) Vliv na výkon | (15:24) Paměť a nové hry | (31:34) Kolik RAM/VRAM mít | (40:04) Aktuální stav trhu | (44:15) Časování a latence}{1:22:07}{39592728}{}
\FC{hardware-club-031.mp3}{\HCtitle{\#31}{Stavíme herní PC (jarní aktualizace)}}{2019-04-25}{\Janchor, \JSvak, \JZubec}{}{(01:46) PC do 10\,000\,Kč | (14:00) Herní PC do 15\,000\,Kč | (31:28) Střední třída do 30\,000\,Kč | (46:15) Dělo za 40\,000\,Kč s RTX 2080 | (55:35) Nejvýkonnější herní PC za 68\,000\,Kč}{1:19:47}{39644100}{}
\FC{hardware-club-032.mp3}{\HCtitle{\#32}{Jubilejní s Michalem Rybkou}}{2019-05-16}{\Janchor, \JSvak, \JZubec, \Wolf}{}{(00:27) Představení hosta | (02:10) Téma: Google Stadia a streamované hraní | (55:40) Várka dotazů | (1:22:06) Soutěž | (1:24:19) Nečekaný vstup}{1:29:28}{42652512}{}
\FC{hardware-club-033.mp3}{\HCtitle{\#33}{Computex 2019}}{2019-05-30}{\Sleepless, \Janchor, \JSvak, \JZubec}{}{(01:05) Nové procesory AMD Ryzen 3000 | (25:41) Ryzeny 3000 a nové konzole | (34:58) Nové grafiky AMD Navi | (53:37) Dotazy (57:35) Procesory Intel 10. generace}{1:03:35}{31435596}{}
% \FC{hardware-club-034.mp3}{\HCtitle{\#34}{Xbox Scarlett, 16jádro AMD a nové grafické karty (#E3 2019)}}{2019-06-14}{\JSvak, \JZubec}{}{(02:19) AMD Ryzen 9 3950X | (22:04) Nové Radeony | (26:44) RX 5700 XT | (37:06) RX 5700 (bez XT) | (45:34) Xbox Scarlett | (1:00:09) Nvidia SUPER}{1:20:18}{43288388}{}
% \FC{hardware-club-035.mp3}{\HCtitle{\#35}{Hraní na Linuxu v roce 2019}}{2019-06-27}{\Janchor, \JSvak, \JZubec}{}{(01:38) Hraní a Linux: oxymóron? | (04:41) Stav Linuxu jako herní platformy | (15:36) Jak a co si lze na Linuxu zahrát | (29:46) SteamPlay/Proton | (43:09) ProtonDB | (46:41) Problémy Linuxu}{1:18:07}{38465424}{}
% \FC{hardware-club-036.mp3}{\HCtitle{\#36}{Výkon procesorů Ryzen 3000 a GPU Navi}}{2019-07-11}{\Janchor, \JSvak, \JZubec}{}{(03:22) Přehled modelů | (15:54) Výkon Ryzen 3000 | (25:22) Přetaktování | (28:51) Single thread | (30:10) Herní testy, vliv RAM | (50:47) Spotřeba | (59:49) Nové grafiky AMD RX 5700 / 5700 XT}{1:24:39}{40874796}{}
% \FC{hardware-club-037.mp3}{\HCtitle{\#37}{Paměti hrozí zdražováním, testy RTX 2080 Super a další novinky}}{2019-07-25}{\Sleepless, \Janchor, \JSvak, \JZubec}{}{(02:22) Situace na trhu s RAM, spor Japonsko vs. Korea | (14:41) Koupit operační paměť teď, nebo čekat? | (20:14) SpaceX | (44:06) Herní AI Google | (46:38) GeForce RTX 2080 Super: stojí za to? | (51:16) Nové Raspberry Pi}{0:59:58}{28244412}{}
% \FC{hardware-club-038.mp3}{\HCtitle{\#38}{Jak si vedou starší procesory v roce 2019?}}{2019-08-15}{\Janchor, \JSvak, \JZubec}{}{(01:22) Výkon procesorů od 486 | (20:24) Vývoj od Sandy Bridge dále | (30:05) Srovnání v Cinebench | (41:05) Herní výkon napříč generacemi | (45:28) GTA V}{1:03:30}{31561344}{}
% \FC{hardware-club-039.mp3}{\HCtitle{\#39}{Nové herní sestavy s Ryzeny a AMD Navi}}{2019-08-29}{\JSvak, \JZubec}{}{1:50 Změny na trhu | 14:53 Základní herní PC za 12000 Kč | 25:39 Počítač za 17000 Kč | 43:31 Herní PC za cca 30000 Kč | 57:26 Sestava High aka co nejvyšší výkon kolem 40000 Kč | 1:06:29 Sestava TOP OC}{1:23:07}{42933024}{}
% \FC{hardware-club-040.mp3}{\HCtitle{\#40}{O nejnáročnějších PC hrách}}{2019-09-12}{\JSvak, \JZubec}{}{5:13 90. léta: DOOM, Descent, Wing Commander, Quake I \&{} II, Total Annihilation, Unreal, Half Life, Outcast | 47:44 Rok 2000+: Giants: Citizen Kabuto, Morrowind, Far Cry, Half Life2, DOOM 3, Fear, TES V: Oblivion | 57:52 Rok 2007 a legendární Crysis | 1:04:07 2008+: GTA IV, S.T.A.L.K.E.R.: Call of Pripyat, Metro Last Light, Crysis 3}{1:25:49}{45336744}{}
% \FC{hardware-club-041.mp3}{\HCtitle{\#41}{O bazarovém hardwaru}}{2019-09-26}{\JSvak, \JZubec}{}{(02:53) Výhody a nevýhody nákupu z bazaru | (10:28) Jaké komponenty z bazaru koupit a jaké raději ne? | (38:52) Kde nakupovat/prodávat + jak se vyhnout podvodníkům | (57:21) Co můžeme dnes najít v bazaru a za kolik? Co se vyplatí? | (1:28:29) Naše zkušenosti, nejlepší/nejhorší koupě z bazaru}{1:34:44}{48602052}{}
% \FC{hardware-club-042.mp3}{\HCtitle{\#42}{O herních klávesnicích}}{2019-10-17}{\JSvak, \JZubec}{}{(02:48) Co je to herní klávesnice a jaké má mít parametry? | (13:38) Podsvícení kláves | (28:51) Mechanické klávesnice| (48:47) Proč jsou mechanické klávesnice hlučné? | (55:18) Exotické a „sci-fi“ klávesnice}{1:27:22}{44062164}{}
% \FC{hardware-club-043.mp3}{\HCtitle{\#43}{O  budoucnosti výroby hardwaru}}{2019-11-01}{\JSvak, \JZubec}{}{(03:21) GeForce GTX 1660 Super | (10:58) Intel Core i9-9900KS | (14:42) Zvýšené TDP a háček v 5GHz boostu | (27:28) Výroba mikročipu – vývoj a dnešní stav | (44:40) EUVL a technologie ASML | (59:33) Budoucnost litografie}{1:25:18}{45489168}{}
% \FC{hardware-club-044.mp3}{\HCtitle{\#44}{Jak měřit FPS ve hrách}}{2019-11-14}{\JSvak, \JZubec}{}{(02:11) Co je to FPS a jak to měřit? | (08:42) Pokročilé možnosti | (22:03) FPS versus frametimes | (38:25) Novinky – faul marketingu Intel + další zranitelnosti CPU | (54:32) AMD Radeon 5500 | (59:54) Drby o Nvidia Ampere}{1:11:58}{40026996}{}
% \FC{hardware-club-045.mp3}{\HCtitle{\#45}{Neznakomp Vašek staví počítač}}{2019-12-05}{\JSvak, \JZubec, \VPechacek}{}{Překoná Václav Pecháček nástrahy konektorů a drátoví?}{1:33:54}{}{}
% \FC{hardware-club-046.mp3}{\HCtitle{\#46}{AMD a Intel na CES 2020 a další best of veletrhu}}{2020-01-09}{\JSvak, \JZubec}{}{Nejzajímavější představení z veletrhu CES 2020 v Las Vegas. Threadripper, APU Renoir a další novinky od AMD a Intel.}{1:18:08}{49921776}{}
% \FC{hardware-club-047.mp3}{\HCtitle{\#47}{Kolik jader procesoru je už moc?}}{2020-01-23}{\JSvak, \JZubec}{}{S vydáním Threadripperu 3990X s 64 jádry / 128 vlákny se na jazyk znova krade otázka, zda je tak rapidní vývoj v počtu jader u spotřebitelských procesorů vůbec k něčemu.}{1:09:28}{40501404}{}
% \FC{hardware-club-048.mp3}{\HCtitle{\#48}{Zpackané vydání Radeonu RX 5600 XT}}{2020-02-06}{\JSvak, \JZubec}{}{Nová grafická karta AMD Radeon RX 5600 XT dostala pod tlakem konkurence na poslední chvíli injekci s vyšším výkonem. Co taková změna způsobila a proč AMD raději kartu nezlevnilo?}{1:00:15}{31151844}{}
% \FC{hardware-club-049.mp3}{\HCtitle{\#49}{Cena PlayStation 5 a dopad koronaviru na HW}}{2020-02-20}{\JSvak, \JZubec}{}{Oproti minulé generaci se na startu prodeje PlayStationu 5 dočkáme zdražení, vyplývá ze zprávy Bloomberg, která odhalila pravděpodobnou výrobní cenu nové konzole. Nemůže za to třeba koronavirus?}{0:56:37}{27834444}{}
% \FC{hardware-club-050.mp3}{\HCtitle{\#50}{O modování her}}{2020-03-05}{\JSvak, \JZubec}{}{Jedni modování odsuzují, pro druhé je to prakticky další náplň hry. Jisté ale je, že mody jsou recept na dlouhověkost her: Skyrim, Half-Life, DOOM, nebyly by bez modů už v propadlišti dějin?}{1:22:55}{43613532}{}
% \FC{hardware-club-051.mp3}{\HCtitle{\#51}{Nové procesory Intel Comet Lake-S: 10/20 jader, ale pořád 14 nm}}{2020-04-30}{\JSvak, \JZubec}{}{Na vydání nových desktopových procesorů Intel se čeká už dlouho. Poslední refresh vyšel před dvěma lety a mezitím AMD zamíchalo kartami. Co novinky přinesou a vrátí Intel zase na trůn PC builderů?}{1:14:05}{33560208}{}
% \FC{hardware-club-052.mp3}{\HCtitle{\#52}{Nové GPU Ampere od Nvidia}}{2020-05-14}{\JSvak, \JZubec}{}{Jensen Huang představil novou architekturu GPU Ampere, ač zatím pouze jako výpočetní čip pro strojové učení. Epic ukázal demo Unreal Engine 5. Oboje si v 52. HWC probereme.}{0:54:40}{26402544}{}
% \FC{hardware-club-053.mp3}{\HCtitle{\#53}{}}{2020--}{\JSvak, \JZubec}{}{}{}{}{}
% \FC{hardware-club-054.mp3}{\HCtitle{\#54}{}}{2020--}{\JSvak, \JZubec}{}{}{}{}{}
% \FC{hardware-club-055.mp3}{\HCtitle{\#55}{}}{2020--}{\JSvak, \JZubec}{}{}{}{}{}
% \FC{hardware-club-056.mp3}{\HCtitle{\#56}{}}{2020--}{\JSvak, \JZubec}{}{}{}{}{}
% \FC{hardware-club-057.mp3}{\HCtitle{\#57}{}}{2020--}{\JSvak, \JZubec}{}{}{}{}{}
% \FC{hardware-club-058.mp3}{\HCtitle{\#58}{}}{2020--}{\JSvak, \JZubec}{}{}{}{}{}
% \FC{hardware-club-059.mp3}{\HCtitle{\#59}{}}{2020--}{\JSvak, \JZubec}{}{}{}{}{}
% \FC{hardware-club-060.mp3}{\HCtitle{\#60}{}}{2020--}{\JSvak, \JZubec}{}{}{}{}{}
% \FC{hardware-club-061.mp3}{\HCtitle{\#61}{}}{2020--}{\JSvak, \JZubec}{}{}{}{}{}
% \FC{hardware-club-062.mp3}{\HCtitle{\#62}{}}{2020--}{\JSvak, \JZubec}{}{}{}{}{}
% \FC{hardware-club-063.mp3}{\HCtitle{\#63}{}}{2020--}{\JSvak, \JZubec}{}{}{}{}{}
% \FC{hardware-club-064.mp3}{\HCtitle{\#64}{}}{2020--}{\JSvak, \JZubec}{}{}{}{}{}
% \FC{hardware-club-065.mp3}{\HCtitle{\#65}{}}{2020--}{\JSvak, \JZubec}{}{}{}{}{}
% \FC{hardware-club-066.mp3}{\HCtitle{\#66}{}}{2020--}{\JSvak, \JZubec}{}{}{}{}{}
% \FC{hardware-club-067.mp3}{\HCtitle{\#67}{}}{2020--}{\JSvak, \JZubec}{}{}{}{}{}
% \FC{hardware-club-068.mp3}{\HCtitle{\#68}{}}{2020--}{\JSvak, \JZubec}{}{}{}{}{}
% \FC{hardware-club-069.mp3}{\HCtitle{\#69}{}}{2020--}{\JSvak, \JZubec}{}{}{}{}{}